\section{Problem definition}\label{sec:problem}

The assignment explores how different algorithmic strategies can analyse, search, compare and
characterise biological sequences. Rather than apply one technique, this report applies every
technique of the course to a concrete question that the problem scenario poses about the
SARS-CoV-2 files, and then compares the techniques on measured cost. The main technique of the
report is Hirschberg's algorithm, dynamic programming for the longest common subsequence
combined with divide and conquer. It is applied to the comparison of the Wuhan 2019 genome
(MN908947.3) with the Texas 2020 genome (MT106054.1). The other techniques are supporting
evidence that answer smaller questions about the same files and provide the baselines against
which the main result is judged.

The mapping between questions and techniques is:
\begin{itemize}\setlength\itemsep{0pt}
    \item Where are the genes M, S and ORF1ab inside the genome? Exact pattern search: brute force as the baseline, then KMP and the Z-function.
    \item How similar are the two genomes and what exactly changed? Longest common subsequence by dynamic programming, made feasible for two full genomes by Hirschberg's divide-and-conquer formulation.
    \item Which open reading frames form the largest non-overlapping set? Greedy activity selection.
    \item Where is each protein of \texttt{seq-proteins.txt} encoded? Backtracking over codon choices with pruning by the genome.
    \item Which motif of length $k$ occurs most often when one mismatch is allowed? Branch and bound over the tree of prefixes.
    \item Which region of each gene is the longest palindrome? Manacher's algorithm against centre expansion (extra technique).
\end{itemize}

Success is measured by three criterion per technique: a correct answer checked against an independent
computation, a justification of why the technique fits the problem, and a measured running time
compared with its theoretical complexity and with a brute-force baseline.
